% !TEX root = ../../eccv2016submission.tex

\section{Related Works}

\noindent\textbf{Text-to-Image Generation.}
Text-to-image (T2I) generation is a conditional generative modeling problem: given a text prompt \(y\), the model samples an image \(x \sim p_\theta(x \mid y)\) that is both visually plausible and semantically aligned with the prompt. T2I lies at the interface of natural language processing and computer vision, and has become a core task in modern AIGC pipelines~\cite{shi2020captioning,ramesh2022hierarchical}.

Recent T2I systems can be grouped into a few practical families. First, hierarchical and autoregressive pipelines model text-conditioned generation via staged decoders (e.g., CLIP-latent priors and image decoders)~\cite{ramesh2022hierarchical}. Second, retrieval-augmented approaches inject external visual neighbors to improve faithfulness and factual grounding~\cite{blattmann2022retrieval,chen2022reimagen}. Third, diffusion-based and flow-based formulations now dominate large-scale generation quality, including latent diffusion variants and transformer backbones~\cite{podell2024sdxl,chen2024pixartsigma}.

In diffusion-style T2I models, the text prompt is encoded into embeddings that condition denoising updates across timesteps. Sampling then iteratively reconstructs an image from noise while preserving prompt semantics. A standard mechanism is classifier-free guidance (CFG), which combines conditional and unconditional predictions to trade off prompt adherence and sample realism~\cite{ho2022classifierfree}. In practice, efficiency-oriented distillation methods can further reduce the required number of denoising steps while retaining competitive quality~\cite{sauer2024add}.

More generally, diffusion models learn generation by reversing a gradual corruption process. A forward Markov process adds noise to real images over many steps, and a learned denoiser approximates the reverse dynamics to map noise back to the data manifold. This iterative formulation is one reason diffusion models are robust for high-dimensional image synthesis: training is stable, conditioning is flexible, and sampling quality scales well with model/data size. Recent systems further improve this paradigm through latent-space modeling and transformer-based rectified-flow formulations~\cite{podell2024sdxl,scaling_rectified_flow_transformers_for_high_resolution_image_synthesis}.

Current state-of-the-art T2I performance is largely achieved by scaled diffusion/flow-transformer systems trained on massive text-image corpora~\cite{podell2024sdxl,chen2024pixartsigma}. Stable Diffusion 3 belongs to this frontier: it uses a rectified-flow transformer formulation and strong text conditioning to improve prompt following, typography, and high-resolution synthesis~\cite{scaling_rectified_flow_transformers_for_high_resolution_image_synthesis}. In this paper, we use SD~3.5 as the semantic composition baseline and compare it against logical score-space composition operators.

\noindent\textbf{Compositional Text-to-Image Generation.}
Compositional text-to-image generation requires a model to satisfy multiple constraints simultaneously in a single image, including object identity, attributes, counts, and inter-object relations. Although modern T2I models produce high-fidelity samples, they still frequently fail on compositional prompts, e.g., by swapping attributes across objects or violating spatial/relational constraints. Prior work has explored compositionality through focused subproblems such as concept conjunction, structured guidance, and attention-based control~\cite{liu2022compositional,kitada-2022-tfsdg,Wu_2023_ICCV}. A complementary evaluation direction is open-world benchmarking: T2I-CompBench shows that compositional robustness remains a key bottleneck even for strong contemporary models~\cite{NEURIPS2023_f8ad010c}. In this paper, we focus on inference-time composition by combining pretrained diffusion experts in score space~\cite{liu2022compositional}.


